% FILE: 01_research_question.tex
% Project: sources-pm4py — PM4Py Comparative Oracle Atlas
% Thesis role: pm4py-comparative-oracle-atlas

\section{Research Question}
\label{sec:sources-pm4py-rq}

\subsection{Problem Statement}
\label{sec:sources-pm4py-rq-problem}

Process mining libraries have long served as reference implementations for conformance
checking, process discovery, and fitness evaluation. PM4Py is the most widely adopted
open-source process mining library in the Python ecosystem. When building a WebAssembly
process intelligence bridge intended to operate in browser and edge environments, a
precise mapping is required: which PM4Py capabilities must be replicated exactly, which
may be approximated, and which would introduce structural defects if ported? This question
cannot be answered by reading API documentation alone. It requires a systematic capability
enumeration, a formal divergence classification, a type-law analysis of what PM4Py's
runtime assumptions would mean if adopted in a Rust/WASM type-enforced context, and a
gap inventory that separates missing type surfaces from missing engine logic. The
\texttt{sources/pm4py} corpus is the research program's answer to this problem: a
complete oracle atlas that defines the boundary between the Python process mining world
and the Rust/WASM process evidence world.

\subsection{Old-Regime Assumption Challenged}
\label{sec:sources-pm4py-rq-oldregime}

The old-regime assumption is that a process mining library is a template: any new
implementation should replicate the existing library's function signatures, data structures,
and runtime patterns. PM4Py accepts \texttt{Union[EventLog, pd.DataFrame]} at nearly every
function boundary, performs runtime \texttt{isinstance} checks, and branches silently into
DataFrame-specific code paths. Discovery results carry no lineage: a Petri net manufactured
by \texttt{pm4py.discover\_petri\_net\_inductive()} carries no record of which log version,
which timestamp, or which algorithm parameters produced it. Validation errors are raised as
generic Python \texttt{Exception} with string messages rather than named law types. Silent
attribute projection loss occurs when \texttt{pm4py.project\_log()} drops columns not in
the requested set. These patterns represent a specific architectural philosophy: maximise
ergonomic flexibility at the cost of type safety and provenance. The research question is
whether this philosophy transfers to a Rust/WASM context, or whether it would destroy the
type-law covenant that makes cryptographic process receipts possible.

\subsection{Hypothesis}
\label{sec:sources-pm4py-rq-hypothesis}

The hypothesis of this corpus is that PM4Py's 120 capabilities divide into three
structurally distinct groups that require different handling in wasm4pm. First, approximately
75 deterministic capabilities (those bridged in \texttt{pm4wasm.d.ts}) require exact
algorithmic equivalence: any divergence would break the audit chain. Second, approximately
40 heuristic or rendering capabilities admit safe divergence as long as output semantics
are preserved and the divergence is documented in \texttt{Evidence} metadata. Third,
ten architectural patterns --- most prominently DataFrame mutability, generic exception
raising, silent projection loss, provenance-free discovery, and LLM chain-of-thought
as a discovery path --- must never be replicated, because they would render the
\texttt{Evidence<T, State, Witness>} type surface meaningless. This tripartite division
is the structural finding of the corpus. [AUTHOR THESIS]

\subsection{Success Criteria}
\label{sec:sources-pm4py-rq-criteria}

Success for the \texttt{sources/pm4py} corpus is defined by three measurable outcomes.
First, every PM4Py function appearing in the WASM bridge (\texttt{pm4wasm.d.ts}) must
have a documented wasm4pm type surface or a GAP status entry with severity classification.
This is certified by the PM4PY\_ORACLE\_RECEIPT in \texttt{receipts/RECEIPT\_REGISTRY.md},
which records 8 CRITICAL/HIGH gaps remaining and tracked in the \texttt{gaps/} directory.
Second, the ALIVE gate criterion 4 must be met: $\geq 4$ files in \texttt{sources/pm4py}
(actual: 14 at assessment, confirmed by \texttt{ALIVE\_GATE\_ASSESSMENT.md}). Third, the
divergence matrix must provide a binary classification for every algorithm --- must-match
or safe-divergence --- with an explicit rationale tied to whether the algorithm produces a
deterministic output that enters the proof chain. The divergence matrix v2.0
(2026-05-31) satisfies this criterion for all 37 assessed algorithms.

\bigskip
\noindent\textit{Source authority:} NOT\_TO\_COPY.md (anti-patterns 1--10),
divergence-matrix.md (sections 11.1--11.3), ALGORITHM\_CROSSWALK.md (Summary Table),
receipts/RECEIPT\_REGISTRY.md (PM4PY\_ORACLE\_RECEIPT), ALIVE\_GATE\_ASSESSMENT.md
(criterion 4).
